\chapter{The Plain Type System}

The plain type system is the foundational language of \DepPy{}. Its home is
\repo{src/deppy/types/}, and its central role is to define the semantic objects that later
passes manipulate. The file \repo{src/deppy/types/base.py} opens by describing itself as a
base hierarchy for \DepPy{}'s sheaf-descent semantic typing system. That phrase is easy to
overlook, but it says something important: the foundational system is already interpreted
semantically and geometrically.

\section*{TypeBase and the carrier language}
Every concrete plain type inherits from \py{TypeBase}. The abstract interface requires
substitution, free-variable computation, structural traversal, equality, hashing, and a
human-readable representation. Those are the operations one needs if types are to be
symbolic objects in a real formal language rather than inert runtime labels.

The concrete hierarchy covers primitive Python types, top and bottom types,
container-oriented types, callables, protocols, classes, and modules. This breadth matters
because it means the plain system aims to be a semantic typing language for native Python,
not a separate toy calculus with only theorem-prover examples.

\section*{Refinement types}
The refinement layer in \repo{src/deppy/types/refinement.py} defines the familiar shape
\[
  \{v : \tau \mid \varphi(v)\}
\]
and then invests serious engineering effort in making predicates behave like real symbolic
objects. Predicates support free-variable analysis, substitution, negation, conjunction,
disjunction, implication, comparison, and named qualifiers. This is the first place where
one sees the repository's long-term ambition: specifications are not meant to be English
comments left floating around; they are intended to be structured semantic objects.

The crucial point for later chapters is that the plain refinement layer has no trust
metadata. A refinement is simply a predicate-bearing type. There is no confidence score,
no provenance, and no solver-versus-oracle distinction internal to the object itself. This
absence is deliberate. It is one of the reasons the hybrid layer later becomes necessary.

\section*{Dependent formers}
The file \repo{src/deppy/types/dependent.py} gives the foundational system four of the
most important ingredients for the whole repository:
\begin{itemize}[leftmargin=2em]
  \item \py{PiType} for dependent functions $\Pi(x : A).B(x)$,
  \item \py{SigmaType} for dependent pairs $\Sigma(x : A).B(x)$,
  \item \py{IdentityType} for equality-like objects,
  \item indexed and quantified families for parameterized semantic structure.
\end{itemize}
These are not decorative imports from dependent type theory. They are what allow the rest
of the repository to speak concretely about APIs whose outputs depend on inputs,
witness-bearing results, transport of equalities, and theorem-carrying program fragments.

\begin{definition}[Plain dependent function]
A plain dependent function type is a triple $(x, A, B)$ represented in code as
\py{PiType(param_name, param_type, return_type)}, with substitution and free-variable
operations respecting binder hygiene.
\end{definition}

\begin{definition}[Plain dependent pair]
A plain dependent pair type is a triple $(x, A, B)$ represented as
\py{SigmaType(fst_name, fst_type, snd_type)}, where the second component type may depend on
the first component.
\end{definition}

\section*{Subtyping and obligations}
The plain system does not stop at representation. The subtype checker in
\repo{src/deppy/types/subtyping.py} produces subtype answers and verification conditions.
This is what turns the type hierarchy into an ordered refinement lattice rather than a mere
syntax tree. When later hybrid code says it values types in a refinement lattice, this is
the layer providing the foundational meaning of that phrase.

\paragraph{Why the plain system still matters after hybridization.}
It is tempting to think the hybrid system replaces the plain one. That reading is wrong.
The mixed-mode compiler still lowers user-facing syntax into plain obligations such as
refinements and Sigma-like witness forms. Lean translation still needs a foundational
type-level language. The agent and checkers still benefit from a symbolic substrate with
substitution and binder structure. The correct statement is therefore:
\begin{quote}
The plain system is the object language; the hybrid system is a richer semantics for how
objects of that language are layered, checked, and trusted.
\end{quote}

\begin{proposition}
If the hybrid system were removed but the plain type system retained, \DepPy{} would still
have a meaningful dependent-and-refinement type language. If the plain system were
removed but the hybrid packaging retained, the repository would lose the exact symbolic
objects that many higher-level workflows require.
\end{proposition}

The proof is architectural rather than formal: the code base repeatedly compiles into,
imports from, or conceptually depends on the plain hierarchy.
